% =============================================================================
%  _borrow_flow_tikz.tex — the borrow-rule flowchart (TikZ body only)
%  Recreated from resource/Flowchart Whiteboard ... Neon Style-1.pdf
%  Requires (loaded in the preamble of whichever document \inputs this):
%     fontspec (Laksaman), xcolor, tikz + libraries
%       shapes.geometric, arrows.meta, positioning, calc, fit, backgrounds
%  Colours are (re)defined here so the file is portable between the proposal
%  and the standalone test wrapper.
% =============================================================================

% ---- light, restrained palette (a calmer take on the neon reference) --------
\definecolor{flStart}{HTML}{CDEFD6}   % green   — start / terminal
\definecolor{flProc} {HTML}{E4EEFB}   % blue    — borrower steps
\definecolor{flStaff}{HTML}{D3EFEA}   % teal    — staff (ผู้ให้ยืม) steps
\definecolor{flDec}  {HTML}{FBF0CC}   % yellow  — decisions
\definecolor{flUse}  {HTML}{FBDFC2}   % orange  — in-use
\definecolor{flEnd}  {HTML}{F6D3E7}   % pink    — final update
\definecolor{flStop} {HTML}{F8CFCF}   % red     — try again
\definecolor{flEdge} {HTML}{5A5A5A}   % node borders / arrows

\begin{tikzpicture}[
    node distance = 6.5mm and 10mm,
    every node/.style = {font=\footnotesize},
    box/.style       = {rectangle, rounded corners=2pt, draw=flEdge, line width=0.4pt,
                        align=center, text width=3.2cm, inner sep=3.5pt, minimum height=8mm},
    startstop/.style = {box, fill=flStart, text width=2.7cm},
    proc/.style      = {box, fill=flProc},
    staff/.style     = {box, fill=flStaff},
    inuse/.style     = {box, fill=flUse, text width=2.7cm},
    endnode/.style   = {box, fill=flEnd, text width=2.9cm},
    stopnode/.style  = {box, fill=flStop, text width=1.9cm, minimum height=7mm},
    decision/.style  = {diamond, aspect=2.1, draw=flEdge, line width=0.4pt, fill=flDec,
                        align=center, inner sep=0pt, text width=2.3cm},
    note/.style      = {rectangle, rounded corners=2pt, draw=flEdge!60, dash pattern=on 2pt off 1.5pt,
                        fill=white, align=left, text width=4.5cm, inner sep=4pt, font=\scriptsize},
    arr/.style       = {-{Stealth[length=2mm]}, draw=flEdge, line width=0.5pt},
    lbl/.style       = {font=\scriptsize, fill=white, inner sep=1pt}
  ]

  % ---------------- main spine (ผู้ยืม / Borrower) --------------------------
  \node[startstop]                     (start)  {เลือกของที่จะยืม};
  \node[proc, below=of start]          (tier)   {จำแนกระดับอุปกรณ์\\T0 \textperiodcentered\ T1 \textperiodcentered\ T2 \textperiodcentered\ T3};
  \node[proc, below=of tier]           (req)    {ส่งคำขอยืม\\(ระบุจำนวน / หมายเลข / ช่วงเวลา ตามระดับ)};
  \node[decision, below=8mm of req]    (auth)   {การอนุมัติ\\ผ่านหรือไม่?};
  \node[staff, below=8mm of auth]      (prep)   {เจ้าหน้าที่จัดเตรียม / เลือกอุปกรณ์\\(T1: ชิ้นพร้อมใช้ \textperiodcentered\ T2: ตามหมายเลข)};
  \node[proc, below=of prep]           (pin)    {ถ่ายรูปตอนรับ + อัปโหลดเข้าระบบ};
  \node[inuse, below=of pin]           (use)    {กำลังใช้งาน\\(IN USE)};
  \node[decision, below=7mm of use]    (renew)  {ต้องการ\\ต่ออายุ?};
  \node[decision, right=14mm of renew, text width=2.2cm] (ronline) {เหลือสิทธิ์\\ต่อออนไลน์?};
  \node[proc,  above=16mm of ronline, text width=2.8cm] (rsys) {ต่อผ่านระบบ (ออนไลน์) 1 ครั้ง};
  \node[staff, below=7mm of ronline, text width=2.8cm] (rchk) {นำอุปกรณ์มาให้เจ้าหน้าที่ตรวจสภาพ};
  \node[proc, below=7mm of renew]      (pout)   {ถ่ายรูปตอนคืน + อัปโหลดเข้าระบบ};
  \node[decision, below=7mm of pout]   (insp)   {ตรวจสอบหลังคืน\\พบความเสียหาย?};
  \node[endnode, below=8mm of insp]    (upd)    {อัปเดตข้อมูล\\ของที่ยืมและผู้ยืม \; (จบ)};

  % ---------------- side branches ------------------------------------------
  \node[stopnode, left=14mm of auth]   (retry)  {Try Again};
  \node[staff, right=12mm of insp, text width=3.0cm] (dmg) {รายงานความเสียหาย\\$\rightarrow$ หักเครดิต (discredit)};

  % tier-specific approval note
  \node[note, right=8mm of auth] (authnote)
    {\textbf{การอนุมัติตามระดับ}\\[1pt]
     T0–T1: ระบบตรวจสิทธิ์ + เครดิต + ของคงเหลือ\\
     T2: อาจารย์ (Supervisor) อนุมัติ\\
     T3: ตรวจช่วงเวลาว่าง (IsFree)};

  % ---------------- arrows: main spine -------------------------------------
  \draw[arr] (start) -- (tier);
  \draw[arr] (tier)  -- (req);
  \draw[arr] (req)   -- (auth);
  \draw[arr] (auth)  -- node[lbl, right]{Accept} (prep);
  \draw[arr] (prep)  -- (pin);
  \draw[arr] (pin)   -- (use);
  \draw[arr] (use)   -- (renew);
  \draw[arr] (renew) -- node[lbl, right]{ไม่} (pout);
  \draw[arr] (pout)  -- (insp);
  \draw[arr] (insp)  -- node[lbl, right]{ไม่พบ} (upd);

  % reject -> try again -> back to request
  \draw[arr] (auth) -- node[lbl, above]{Reject} (retry);
  \draw[arr] (retry.north) |- (req.west);

  % renewal cycle: 1 online renewal, then a staff condition check, alternating
  % (a passed check re-earns the next online renewal — see caption). Both
  % outcomes MERGE onto a right-hand rail and return to IN USE via a single
  % arrow, so no return line crosses the online / staff-check nodes.
  \draw[arr] (renew.east)    -- node[lbl, above]{ใช่} (ronline.west);
  \draw[arr] (ronline.north) -- node[lbl, right]{มี} (rsys.south);
  \draw[arr] (ronline.south) -- node[lbl, right]{ไม่มี} (rchk.north);
  \coordinate (rrx) at ($(rsys.east)+(0.8,0)$);
  \draw (rsys.east) -- (rsys.east -| rrx);
  \draw (rchk.east) -- (rchk.east -| rrx);
  \draw (rsys.east -| rrx) -- (rchk.east -| rrx);
  \draw[arr] (use.east -| rrx) -- (use.east);

  % damage found -> staff report -> converge to update
  \draw[arr] (insp) -- node[lbl, above]{พบ} (dmg);
  \draw[arr] (dmg.south) |- (upd.east);

  % dashed link to the approval note
  \draw[flEdge!60, dash pattern=on 2pt off 1.5pt] (auth.east) -- (authnote.west);

\end{tikzpicture}
